\section{Integral and $\lambda$-ring controls on the bridges}\label{sec:root-findings}
The results in this section are deductions checked in this run. They are not claims of historical priority. We distinguish an identity of characters, a decomposition in a Grothendieck group, an actual direct-sum decomposition, and a dimension identity.

\begin{proposition}[G-01: The three meanings of the second operation]\label{G-01}
Let $E$ be a finite $G$-set. Write $b_2(E)$ for the $G$-set of two-element subsets and $s_2(E)$ for the $G$-set of two-element multisets. In the Burnside ring one has
\[
 [s_2(E)]=[b_2(E)]+[E],\qquad [E]^2=[E]+[\operatorname{Inj}(\{1,2\},E)].
\]
The identity $[E]^2=2[b_2(E)]+[E]$ does not hold in general. If $V$ is a representation over a field of characteristic zero, the correct special $\lambda$-ring identities are
\[
 [\operatorname{Sym}^2V]=[V]^2-[\wedge^2V],\quad
 [\operatorname{Sym}^2V]-[\wedge^2V]=\psi^2[V],\quad
 [V]^2=2[\wedge^2V]+\psi^2[V].
\]
Consequently the expression $2\lambda^2(E)+E=E^2$ in cote 133, scan 019, line 53 is valid after taking cardinalities, but is not a universal equivariant or special-$\lambda$-ring identity. The supplied study's assertion that the two equivariant descriptions agree needs this correction.
\end{proposition}
\begin{proof}
Partition a multiset of cardinality two into the diagonal multisets $\{e,e\}$ and the subsets $\{e,e'\}$ with $e\ne e'$. This partition is $G$-equivariant. Similarly partition ordered pairs into diagonal and off-diagonal pairs. For a counterexample take $G=C_2$ and its regular two-element set $E$. Then $E\times E$ is a disjoint union of two free $C_2$-orbits, whereas $b_2(E)$ is one fixed point. Thus the asserted equality would read $2[E]=[E]+2[\mathrm{pt}]$. Applying the fixed-point homomorphism for the whole group gives $0=2$, a contradiction.
For representations, first prove the universal identity on $\mathrm{GL}_r$ and then pull it back along the representation of $G$. Write the character of its defining representation on the diagonal torus as $x_1+\cdots+x_r$. The symmetric and exterior squares have characters $\sum_{i\leq j}x_ix_j$ and $\sum_{i<j}x_ix_j$. Their sum is $(\sum_i x_i)^2$, and their difference is $\sum_i x_i^2=\psi^2[V]$. Character injectivity in characteristic zero proves the formulas. In the same $C_2$ example the permutation representation is $1\oplus\epsilon$, and $\wedge^2V=\epsilon$, whereas linearization of $b_2(E)$ is $1$. This explicitly shows why the subset operation cannot be silently identified with exterior power. Taking dimensions sends $\psi^2[V]$ to $\dim V$ and recovers the numerical identity. None of this determines the ink on a scan which has not been supplied.
\end{proof}

\begin{proposition}[G-02: The terminal exterior-power variable is necessary]\label{G-02}
For $i\geq2$, a universal formula for $\lambda^i(xy)$ in special $\lambda$-rings cannot in general omit $\lambda^i(x)$ or $\lambda^i(y)$. Thus the terminal index in cote 133, scan 006, lines 25--26 must be $i$ if those lines state the universal product formula.
\end{proposition}
\begin{proof}
Set $y=1$, for which $\lambda^1(1)=1$ and $\lambda^j(1)=0$ for $j>1$. A formula using only $x,\lambda^2(x),\ldots,\lambda^{i-1}(x)$ would express $\lambda^i(x)$ in the subring generated by these earlier operations. Evaluate at the defining representation of $\mathrm{GL}_i$ and restrict to its diagonal torus. This would express the elementary symmetric polynomial $e_i$ in $\mathbb Z[e_1,\ldots,e_{i-1}]$, contrary to the algebraic independence of $e_1,\ldots,e_i$ (equivalently the fundamental theorem of symmetric polynomials). Interchange $x$ and $y$ for the other endpoint. The argument establishes the required mathematical reading, not a new visual transcription.
\end{proof}

\begin{proposition}[G-03: Finite $\gamma$-data do not determine multiplicities]\label{G-03}
For every integer $D\geq0$ there are two honest representations $A_D,B_D$ of the split group scheme $\mathbb G_{m,\mathbb Z}$, each of rank $2^D$, which have the same class modulo $F^{D+1}_\gamma$ but different weight multiplicities. Over $\mathbb Q$ their Chern characters agree in degrees $0,\ldots,D$.
\end{proposition}
\begin{proof}
Let $L$ denote the standard character and put $r=D+1$. Define
\[
 A_D=\bigoplus_{\substack{0\leq j\leq r\\j\text{ even}}}L^{\otimes j\,\oplus\binom rj},\qquad
 B_D=\bigoplus_{\substack{0\leq j\leq r\\j\text{ odd}}}L^{\otimes j\,\oplus\binom rj}.
\]
The two binomial parity sums are $2^{r-1}$. In the representation ring $\mathbb Z[L,L^{-1}]$ the difference of the classes is $(1-L)^r$, a product of $r$ rank-zero elements and hence an element of $F^r_\gamma$. The constant weight has multiplicity one in $A_D$ and zero in $B_D$, so the characters differ. With $c=c_1(L)$, the difference of the Chern characters is $(1-e^c)^r$, whose lowest possible degree is $r$. This proves the assertion. For a line bundle on a scheme, applying the associated-bundle map gives the identical statement in its $K_0$. This does not assert that the two resulting bundles are nonisomorphic on every scheme.
\end{proof}

\begin{proposition}[G-04: An integral coefficient model and its exact specialization gate]\label{G-04}
Let $W_{\mathbb Z}$ be a finite free module, let $\rho:\mathrm{GL}_{N,\mathbb Z}\to\mathrm{GL}(W_{\mathbb Z})$ be a homogeneous polynomial representation of positive degree $e$, and let $f\in W_{\mathbb Z}$. Define $X_{\mathbb Z}$ as the affine schematic closure of its orbit. Its coordinate algebra is torsion-free, hence flat over $\mathbb Z$. In degree $d$ its orbit-coefficient map is an explicit finite integer matrix
\[
 A_d:\operatorname{Sym}^d(W_{\mathbb Z}^*)\longrightarrow
 \mathbb Z[g_{11},\ldots,g_{NN}]_{\leq ed},\qquad P\longmapsto P(\rho(g)f).
\]
Its rank over $\mathbb Q$ is the degree-$d$ coordinate-ring dimension. If $\delta_d$ is any nonzero maximal-rank minor of $A_d$, then this dimension and the degree-$d$ kernel commute with base change over $\mathbb Z[1/\delta_d]$. At a prime dividing $\delta_d$, rank must be tested anew; the fiber of the integral closure need not be the schematic orbit closure of the reduced point.
\end{proposition}
\begin{proof}
The coordinate homomorphism of the orbit map has target $\mathbb Z[\mathrm{GL}_N]=\mathbb Z[g_{ij},\det(g)^{-1}]$, a torsion-free $\mathbb Z$-module. Its kernel $I$ defines the schematic closure and its image $\mathbb Z[W]/I$ is torsion-free. A module over the principal ideal domain $\mathbb Z$ is flat if it is torsion-free, proving the first assertion. Since the representation is polynomial, each degree-$d$ homogeneous coordinate polynomial maps to a polynomial of degree at most $ed$ in the entries of $g$. Equality on invertible matrices is equality in the polynomial ring because localization by the nonzero polynomial $\det(g)$ is injective. This proves the finite coefficient-matrix description.
Let $r=\operatorname{rank}_{\mathbb Q}A_d$. All $(r+1)$-minors vanish as integers, while an $r$-minor $\delta_d$ is nonzero. Over $\mathbb Z[1/\delta_d]$, use that invertible block for row and column elimination. The vanishing of larger minors makes the remaining Schur complement zero, so the matrix is equivalent to $\operatorname{diag}(I_r,0)$. Its image and kernel are direct summands there and commute with every base change. Without inversion, an inclusion of lattices need not remain injective modulo a prime: already multiplication by $p$ on $\mathbb Z$ has rank one over $\mathbb Q$ and rank zero modulo $p$. Thus no unsupported all-primes orbit-fiber conclusion follows. Homogeneity sends degree $d$ to degree $ed$, so the kernel is homogeneous and the stated grading is well defined. If $r=0$, take $\delta_d=1$.
\end{proof}

\begin{proposition}[G-05: Ternary saturation and the three-row exclusion]\label{G-05}
Work over $\mathbb C$, and let $n\geq2$, $N=n^2$, and $\mathrm{Det}_n\subseteq\operatorname{Sym}^n(\mathbb C^N)^*$ be the determinant orbit closure. Then $\mathrm{Det}_n$ contains every degree-$n$ form supported on a fixed three-dimensional subspace. Consequently, for every partition $\lambda\vdash nd$ with $\ell(\lambda)\leq3$,
\[
 \operatorname{mult}_{S_\lambda\mathbb C^N}\mathbb C[\mathrm{Det}_n]_d
 =\operatorname{mult}_{S_\lambda\mathbb C^N}\operatorname{Sym}^d(\operatorname{Sym}^n\mathbb C^N).
\]
No multiplicity obstruction with at most three rows can separate any other orbit closure of a degree-$n$ form in the same ambient space, including a padded permanent, from $\mathrm{Det}_n$.
\end{proposition}
\begin{proof}
We use the established linear determinantal representation theorem for smooth plane curves: a smooth plane degree-$n$ curve together with a line bundle $L$ of degree $g-1$ with $H^0(L)=0$ admits an $n\times n$ matrix of linear forms whose determinant is its defining form up to a nonzero scalar (Beauville, 2000, Proposition 3.1(a)). Such an $L$ exists over $\mathbb C$: for $g\geq1$, the effective locus in $\operatorname{Pic}^{g-1}$ is the image of the symmetric power $C^{(g-1)}$, of dimension at most $g-1$, whereas $\operatorname{Pic}^{g-1}$ has dimension $g$; for $g=0$, take the degree-$-1$ line bundle. Scaling one row absorbs the nonzero scalar. Hence every smooth ternary degree-$n$ form is obtained by a linear substitution in $\det_n$. Any endomorphism substitution is a limit of invertible substitutions because $\mathrm{GL}_N$ is Zariski dense in $\operatorname{End}(\mathbb C^N)$; continuity of the polynomial action therefore puts these forms in $\mathrm{Det}_n$. Smooth ternary forms form a dense open set, so its closedness puts every ternary form there.
Let $W=\operatorname{Sym}^d(\operatorname{Sym}^n\mathbb C^N)$, and choose the Borel convention whose highest-weight vector has nonnegative weight $\lambda=(\lambda_1,\lambda_2,\lambda_3,0,\ldots,0)$. In a monomial basis of $W$, zero weight in positions $4,\ldots,N$ forces every coefficient variable to involve only the first three basis vectors. Thus the entire highest-weight space $W_\lambda^{\mathrm{hw}}$ is contained in $\operatorname{Sym}^d(\operatorname{Sym}^n\mathbb C^3)$. A nonzero polynomial in these coefficient variables cannot vanish on all ternary forms. By the first paragraph, the degree-$d$ determinant ideal therefore meets $W_\lambda^{\mathrm{hw}}$ trivially. Exactness and complete reducibility in characteristic zero give the asserted multiplicity equality. Every other closed subvariety has a degree-$d$ coordinate ring which is a quotient of $W$, so its multiplicity cannot exceed this common ambient value. This proves the exclusion. The assertion is about complex characteristic-zero multiplicities, not an unqualified statement in positive characteristic.
\end{proof}

\begin{proposition}[G-06: The four-variable gate is nonvacuous]\label{G-06}
For $n\geq4$, the closure $\mathcal D_{n,4}$ of degree-$n$ quaternary forms which are determinants of $n\times n$ matrices of linear forms has dimension at most $2n^2+2$. Since
\[
 \binom{n+3}{3}-(2n^2+2)=\binom{n-1}{3}>0,
\]
this is a proper subvariety of the space of quaternary forms. The determinant ideal in $\operatorname{Sym}(\operatorname{Sym}^n\mathbb C^{n^2})$ therefore has a nonzero isotypic component with a partition of exactly four rows in some finite degree.
\end{proposition}
\begin{proof}
Consider the polynomial map
\[
 (A_1,A_2,A_3,A_4)\longmapsto\det(x_1A_1+x_2A_2+x_3A_3+x_4A_4)
\]
from an affine space of dimension $4n^2$. The group
$H=\{(P,Q)\in\mathrm{GL}_n^2:\det(P)=\det(Q)\}$ acts by $A_i\mapsto PA_iQ^{-1}$, preserving this map. It has dimension $2n^2-1$. Its generic stabilizer has dimension one. Indeed, for a tuple with $A_1=I$, a diagonal $A_2$ with pairwise distinct eigenvalues, and an $A_3$ with every off-diagonal entry nonzero, a stabilizer has $Q=P$; commuting with $A_2$ makes $P$ diagonal, and commuting with $A_3$ makes it scalar. The scalar pairs $(cI,cI)$ always stabilize every tuple. Upper semicontinuity of stabilizer dimension, or the rank calculation for the commutation equations on this exhibited open set, gives generic dimension one. Thus generic fibers contain orbits of dimension $2n^2-2$, and the dimension-of-fibers theorem yields the stated image bound. The displayed identity follows by expanding the binomial coefficient.
Let $\pi$ set all but the first four variables to zero. For every form in $\mathrm{Det}_n$, its image under $\pi$ lies in $\mathcal D_{n,4}$: approximate the form by invertible substitutions of $\det_n$, apply $\pi$, and use the closedness of $\mathcal D_{n,4}$. Choose a nonzero homogeneous equation $P$ of this proper cone. Then $P\circ\pi$ is a nonzero homogeneous equation of $\mathrm{Det}_n$, using only first-four coefficient variables. The ideal of $\mathcal D_{n,4}$ is $\mathrm{GL}_4$-stable; choose a highest-weight vector in a nonzero irreducible submodule of its degree-$d$ part. Embedded in the full coordinate ring, it is also a $\mathrm{GL}_{n^2}$ highest-weight vector: raising operators involving a variable beyond the fourth annihilate it. Its highest weight has at most four rows, and G-05 rules out at most three. Hence its length is exactly four. Existence of such an equation does not establish that it is nonzero on the padded permanent.
\end{proof}

\begin{proposition}[G-07: An exact four-row determinant-closure formula]\label{G-07}
Set $D=320112$. For every $d\geq0$ and partition $\lambda\vdash4d$ with $\ell(\lambda)\leq4$, padded by zeros to length four, the degree-$d$ determinant-closure multiplicity satisfies
\[
 m_{\mathrm{Det}_4}(\lambda,d)
 =a_\lambda(d[4])-a_{\lambda-(D,D,D,D)}((d-D)[4]).
\]
The second term is defined to be zero if $d<D$ or $\lambda_4<D$. In particular the four-row sector is ambient for $d<D$, and at $d=D$ the multiplicity of $\lambda=(D^4)$ is exactly one less than its ambient plethysm multiplicity.
\end{proposition}
\begin{proof}
Leal, Lozano Huerta, and Vite (2024), Theorem 2 and Table 2, identify the closure of linear determinantal quartic surfaces as an irreducible reduced hypersurface of degree $D$ in $\mathbb P(\operatorname{Sym}^4\mathbb C^4)^*$. Table 2 identifies this component $F_1$ by $a=(5,5,5,5)$, $b=(6,6,6,6)$, so every matrix entry has degree one. Its affine cone has a prime principal ideal $(F)$ in the polynomial ring $A=\operatorname{Sym}(\operatorname{Sym}^4\mathbb C^4)$, with $F$ homogeneous of degree $D$.
The line $\mathbb CF$ is stable under $\mathrm{GL}_4$. Every algebraic one-dimensional character of $\mathrm{GL}_4$ is a power of the determinant. Since scalar $zI$ acts on $A_D$ by $z^{4D}$, this character is $\det^D$. Thus multiplication by $F$ gives the exact sequence of $\mathrm{GL}_4$-modules
\[
 0\longrightarrow \det^D\otimes A_{d-D}\xrightarrow{\ \cdot F\ }A_d
 \longrightarrow(A/(F))_d\longrightarrow0.
\]
For $d<D$ its first term is zero.
Let $\pi$ restrict forms in sixteen variables to their first four variables. The image $\pi(\mathrm{Det}_4)$ is precisely the affine cone just described. One inclusion follows by continuity from the orbit of $\det_4$. For the converse, every four-variable linear determinant is an endomorphism substitution in $\det_4$, hence belongs to its orbit closure; the same is true of every limit of such forms, and such a form is fixed by $\pi$. It follows that the intersection of the full determinant ideal with the first-four coefficient polynomial ring is exactly $(F)$. Highest-weight spaces for partitions of at most four rows use only these coefficient variables, by the weight argument of G-05. Applying complete reducibility to the exact sequence and shifting a highest weight by the determinant character proves the formula. At $d=D$ the source is $\det^D\otimes\mathbb C$, so the deficit is one.
\end{proof}
\paragraph{Source correction for G-07.}
The attached Landsberg source, section \texttt{zhanglee}, theorem \texttt{zlthm}, records $640224$. The later result used here is the degree of the reduced prime divisor. Leal--Lozano Huerta--Vite, section 4.1 immediately before Theorem 2, explicitly gives $P_{20,2}=\tfrac12D_{23,14}$. This is the precise cycle-multiplicity distinction; using the unreduced Noether--Lefschetz coefficient as the polynomial degree doubles the answer. This correction concerns the cited literature, not a purported visual correction of cote 133.

\begin{proposition}[G-08: A uniform low-variable determinantal construction]\label{G-08}
Let $m\geq1$ and $k\geq2$, put $a=\lfloor k/2\rfloor$, $b=\lceil k/2\rceil$, and
\[
 B(m,k)=\binom{m+a}{a}+\binom{m+b}{b}-1.
\]
Over any commutative ring, every degree-$m$ homogeneous polynomial $q$ in $k$ variables admits a matrix of linear forms of size $B(m,k)$ with determinant
\[
 \ell^{B(m,k)-m}q,
\]
where $\ell$ may first be an independent variable and then be specialized to any linear form. The construction is polynomial in the coefficients of $q$ and defined over $\mathbb Z$.
Consequently, over $\mathbb C$, whenever $n\geq B(m,k)$, no multiplicity obstruction of length at most $k$ can separate the orbit closure of an arbitrarily padded degree-$m$ form $\ell^{n-m}f$ from $\mathrm{Det}_n$.
\end{proposition}
\begin{proof}
Split the variables into disjoint lists $x=(x_1,\ldots,x_a)$ and $y=(y_1,\ldots,y_b)$, and write
\[
 q=\sum_{|\alpha|+|\beta|=m}c_{\alpha\beta}x^\alpha y^\beta.
\]
Make a rooted directed tree with one vertex for every $x$-monomial of degree at most $m$. The parent of a nonconstant monomial is obtained by deleting one occurrence of its largest-index variable; orient the edge away from the constant root and label it by that variable. There is exactly one root-to-vertex path, with label product $x^\alpha$. Make the corresponding $y$-tree with every edge reversed, so the unique path from vertex $\beta$ to its constant root has product $y^\beta$. Connect vertex $\alpha$ in the first tree to vertex $\beta$ in the second by an edge of constant label $c_{\alpha\beta}$ exactly when $|\alpha|+|\beta|=m$. Every source-to-sink path uses exactly one connecting edge; its label product is $c_{\alpha\beta}x^\alpha y^\beta$. Therefore its path polynomial is exactly $q$. The graph is acyclic and has
\[
 s=\binom{m+a}{a}+\binom{m+b}{b}=B(m,k)+1
\]
vertices.
Choose a topological ordering with source first and sink last, and let $A$ be its strictly upper triangular weighted adjacency matrix. Then $\det(I-A)=1$ and $(I-A)^{-1}=I+A+\cdots+A^{s-1}$. The $(1,s)$ entry is the source-to-sink path polynomial $q$. The adjugate identity therefore expresses $q$, up to the fixed sign $(-1)^{s+1}$, as the determinant of the $(s-1)\times(s-1)$ minor of $I-A$ obtained by deleting row $s$ and column $1$. Change the sign of one row if needed. Its entries are affine linear in $x,y$.
Replace each constant entry $c$ of this minor by $c\ell$. Its determinant is the degree-$B(m,k)$ homogenization of $q$, hence is $\ell^{B(m,k)-m}q$. This identity is universal over $\mathbb Z[c_{\alpha\beta},x,y,\ell]$. Padding the matrix with $n-B(m,k)$ diagonal entries $\ell$ proves the corresponding size-$n$ statement.
For the multiplicity assertion, restrict any invertible substitution of $\ell^{n-m}f$ to $k$ coordinate variables. The result has the form $L^{n-m}q$, with $q$ a homogeneous degree-$m$ form in these variables. The construction puts this result in $\mathrm{Det}_n$. Continuity gives the same inclusion for the entire padded orbit closure. If $P$ is any highest-weight polynomial of length at most $k$ vanishing on $\mathrm{Det}_n$, it uses only coefficient variables supported on the first $k$ underlying variables (in the sense made precise in G-05), and therefore vanishes on the other orbit closure as well. Thus the highest-weight space of the determinant ideal embeds in that of the padded-form ideal. Subtracting their dimensions from the common ambient highest-weight multiplicity gives $m_{\mathrm{padded}}(\lambda,d)\leq m_{\mathrm{Det}_n}(\lambda,d)$ in every degree.
\end{proof}
\paragraph{Quantified implication.}
For $k=4$, the construction gives $B(m,4)=(m+1)(m+2)-1$. More generally $B(m,k)=O_k(m^{\lceil k/2\rceil})$. Thus for a fixed number of rows this sector is excluded once the padding exceeds the displayed bound. This is an explicit restricted lower-bound obstruction theorem, not a general separation result and not an efficient solution of arbitrary plethysm.

\begin{proposition}[G-09: A sharper low-row exclusion by shared exponent chains]\label{G-09}
Let $m\geq1$ and $k\geq3$. Define the following integer bounds. If $k=2r+1$ is odd, set
\[
 C(m,k)=2\binom{m+r}{r}-1.
\]
If $k=2r+2\geq4$ is even, for $1\leq h\leq m+1$ put $Q=\lfloor m/h\rfloor$, $s_0=\min(h-1,m)$ and
\[
\begin{aligned}
 C_h(m,k)&=\sum_{q=0}^{Q}\binom{m-qh+r}{r}+Q(h-1)
          +\sum_{s=0}^{s_0}\binom{m-s+r}{r}-1,\\
 C(m,k)&=\min_{1\leq h\leq m+1}C_h(m,k).
\end{aligned}
\]
Over a field, for every nonzero linear form $L$ and homogeneous degree-$m$ form $q$ in $k$ variables, $L^{C(m,k)-m}q$ has a linear determinantal representation of size $C(m,k)$. In coordinates with $L=x_0$ the construction is defined over $\mathbb Z$ and polynomial in the coefficients of $q$. For fixed $k\geq3$,
\[
 C(m,k)=O_k\bigl(m^{(k-1)/2}\bigr).
\]
Consequently, over $\mathbb C$, if $n\geq C(m,k)$, there is no multiplicity obstruction with at most $k$ rows separating any padded degree-$m$ form from $\mathrm{Det}_n$. This strengthens G-08 in the specified range.
\end{proposition}
\paragraph{Prior-result boundary.}
The asymptotic determinantal-representation order used here is established: Boralevi, van Doornmalen, Draisma, Hochstenbach, and Plestenjak (2017), Theorem 1.3, prove order $m^{v/2}$ for a fixed number $v\geq2$ of affine variables, with a matching generic lower bound. Their proof in section 4 uses connected monomial spaces and an even/odd binary-digit split. We do not claim that order as new. The proof below gives the displayed finite integer bound by a different shared-chain construction, records its coefficient-ring scope, and proves the resulting GCT multiplicity exclusion. These are the forms needed here that the cited theorem does not state.
\begin{proof}
Extend $L$ to a basis of linear forms and write it as $x_0$. Dehomogenize $q$ by setting $x_0=1$, obtaining an arbitrary polynomial $p$ of degree at most $m$ in $k-1$ variables. It suffices to construct an affine linear determinant of size $C(m,k)$ equal to $p$. Homogenizing each matrix entry by $x_0$ then gives $x_0^{C(m,k)-m}q$.

For odd $k=2r+1$, split the dehomogenized variables into lists $x$ and $y$ of $r$ variables each. Use the two monomial trees from G-08, each with $\binom{m+r}{r}$ vertices, and connect every pair $(\alpha,\beta)$ with $|\alpha|+|\beta|\leq m$ by its coefficient in $p$. Every source-to-sink path is unique for its two monomial indices, so its path polynomial is $p$. The cofactor construction in G-08 has size one less than the total number of vertices, giving the displayed bound.

For even $k=2r+2$, write the variables as $x_1,\ldots,x_r,y_1,\ldots,y_r,z$ and fix $h$. Decompose the exponent of $z$ uniquely as $e=qh+s$, $0\leq s<h$. On the left, make a directed chain with $Qh$ edges, all labelled $z$, beginning at the source. At its vertex of index $qh$ attach the outward $x$-monomial tree of degrees at most $m-qh$, identifying its constant root with that chain vertex. Every such tree vertex indexed by $(q,\alpha)$ has a unique source path, with product $z^{qh}x^\alpha$. Chain vertices whose indices are not multiples of $h$ are transit vertices and receive no connecting edges. The left graph has
\[
 (Qh+1)+\sum_{q=0}^Q\left[\binom{m-qh+r}{r}-1\right]
 =\sum_{q=0}^Q\binom{m-qh+r}{r}+Q(h-1)
\]
vertices.

On the right, make a chain with vertices $0,\ldots,s_0$, directed toward vertex $0$, the sink, and again label its edges by $z$. At vertex $s$ attach the inward $y$-monomial tree of degrees at most $m-s$, identifying the root with that vertex. Every tree vertex indexed by $(s,\beta)$ has a unique sink path with product $y^\beta z^s$. The right graph has $\sum_{s=0}^{s_0}\binom{m-s+r}{r}$ vertices. Connect left vertex $(q,\alpha)$ to right vertex $(s,\beta)$ with the constant coefficient of $x^\alpha y^\beta z^{qh+s}$ in $p$ when $|\alpha|+|\beta|+qh+s\leq m$.

All connecting edges go from left to right. Thus every source-to-sink path uses exactly one such edge, and the uniqueness of the two tree paths and the division $e=qh+s$ gives exactly the required monomial with its coefficient, once. The path polynomial is $p$. Taking the same cofactor of $I-A$ as in G-08 yields an affine linear determinant of size $C_h(m,k)$. This graph has a path capable of realizing $z^m$ and a connecting edge, so its number of vertices is at least $m+2$; hence $C_h(m,k)\geq m+1$. The odd bound is also greater than $m$. Homogenization therefore uses a nonnegative exponent. Minimizing over the finite choices of $h$ proves the stated representation bound.

For the asymptotic statement in the even case take $h=\lceil\sqrt m\rceil$. The number of binomial summands is at most $m/h+1+h\leq2\sqrt m+2$, each bounded by $\binom{m+r}{r}=O_r(m^r)$, while $Q(h-1)\leq m$. Thus $C(m,k)=O_r(m^{r+1/2})$. The odd case is immediate from its binomial formula. These exponents are $(k-1)/2$ in both cases.

For the GCT consequence restrict a substituted padded form to any fixed $k$ underlying variables. If its restricted padding $L$ is nonzero, apply the determinant construction and pad with $n-C(m,k)$ diagonal copies of $L$. If $L=0$, the form is zero because $n\geq C(m,k)>m$. Thus every restriction lies in $\mathrm{Det}_n$. Apply the highest-weight ideal argument from G-08, which uses exactly this containment, to obtain the multiplicity inequality in every degree.
\end{proof}
\paragraph{Integral scope and growth of the remaining gate.}
With $L=x_0$ the graph and cofactor identity are universal over the coefficient ring. Extending an arbitrary nonzero $L$ to coordinates is a field operation; over a general base one instead works on a chart where a coefficient of $L$ is invertible and chooses the resulting local basis. No global polynomial choice of that basis is asserted. If $n=m^c$ with fixed $c>0$, every fixed $k\geq3$ satisfying $(k-1)/2<c$ is excluded for sufficiently large $m$. This is a quantitative restriction on possible multiplicity obstructions, not a claim that the remaining higher-row obstruction exists.
\begin{proposition}[G-10: An explicit degree bound for the four-row dimension gaps]\label{G-10}
For $n\geq4$ put
\[
 r_n=2n^2+2,\qquad B_n=n^{r_n-1},\qquad
 d_n=r_n(B_n-1)+1.
\]
For every integer $d\geq d_n$, some partition $\lambda\vdash nd$
with exactly four rows satisfies
\[
 m_\lambda(\operatorname{per}_n)>m_\lambda(\det_n).
\]
More generally, the same explicit threshold $d_n$ applies to every
padded pair $(m,n)$ for which the dimension criterion
$\dim\mathcal P_{m,4}+3>r_n$ of R-14 has been established.
The degree bound is very large and does not specify which partition
has the gap. It makes the preceding ``sufficiently large degree''
assertions effective; it gives no larger range of padded matrix sizes.
\end{proposition}
\begin{proof}
Write $Y=\mathbb P(\mathcal D_{n,4})$, $s=\dim Y$, and
$\Delta=\deg Y$. The pencil coefficient map projectivizes to a
rational map
\[
 \phi:\mathbb P^{4n^2-1}\dashrightarrow Y
\]
given by homogeneous polynomials of degree $n$. Its image is dense.
We first check the degree estimate
\begin{equation}\label{G10:image-degree}
 \Delta\leq n^s\leq n^{r_n-1}=B_n.
\end{equation}
Choose $s$ general hyperplanes in the target. Their intersection
with $Y$ consists of $\Delta$ distinct points in a dense open subset
contained in the actual image of $\phi$. Such an open subset exists
because the dominant image is constructible. Each point therefore
has a nonempty fiber in the domain of definition $U$ of $\phi$.
The common zero locus $Z$ of the $s$ pulled-back degree-$n$
equations in $\mathbb P^{4n^2-1}$ meets $U$ in the disjoint union
of these fibers. Every irreducible component of $Z$ meeting $U$
maps to just one of the target points. At least $\Delta$ distinct
components of $Z$ are consequently required. Each has degree at
least one, while repeated projective B\'ezout gives that the sum
of degrees of all components of $Z$ is at most $n^s$.
This proves the first inequality. Components contained in the
base locus only contribute to that upper bound and cannot cover
any of the nonempty fibers in $U$.
The second inequality follows from G-06, which gives $s\leq r_n-1$.
This is the usual image-degree argument with its base locus
accounted for explicitly; compare Sombra (1997), section~1,
printed preprint p.~5, and Proposition~2.10, pp.~11--12,
whose displayed statement treats an everywhere regular map.

The all-degree Hilbert bound of Chardin (1989), the unnumbered
central theorem on printed p.~306, applies to the homogeneous
prime ideal of the reduced irreducible variety $Y$ over
$\mathbb C$. Thus for every $d\geq1$,
\begin{equation}\label{G10:det-Hilbert}
 \dim\mathbb C[\mathcal D_{n,4}]_d
 \leq\Delta\binom{d+s}{s}
 \leq B_n\binom{d+r_n-1}{r_n-1}.
\end{equation}
We use this cited bound rather than assuming that the determinant
coordinate ring is Cohen--Macaulay.

Let $X$ be either $\mathcal P_{n,4}$ or a padded image
$\mathcal Z_{m,n}$ satisfying the stated criterion. It is a
reduced irreducible standard graded affine cone and has dimension
at least $r_n+1$, by R-15 or R-14 respectively. Linear Noether
normalization over the infinite field $\mathbb C$ supplies
$\dim X$ algebraically independent linear coordinate functions.
Any $r_n+1$ of these generate a polynomial subalgebra of its
coordinate ring. Counting its degree-$d$ monomials gives
\begin{equation}\label{G10:per-Hilbert}
 \dim\mathbb C[X]_d\geq\binom{d+r_n}{r_n}.
\end{equation}
The ratio of this binomial to the one in
\eqref{G10:det-Hilbert} is $(d+r_n)/r_n$.
When $d\geq r_n(B_n-1)+1$, that ratio is strictly greater
than $B_n$. Hence the Hilbert function of $X$ is strictly larger
in every such degree.

Decomposing the two degree-$d$ coordinate modules into
$\mathrm{GL}_4$ irreducibles, their dimension difference is
\[
 \sum_{\lambda\vdash nd,\,\ell(\lambda)\leq4}
 \bigl(m_\lambda(X)-m_\lambda(\mathcal D_{n,4})\bigr)
 \dim S_\lambda\mathbb C^4>0.
\]
At least one multiplicity difference is positive. The
highest-weight restriction argument of R-14 identifies these
multiplicities with those of the original orbit closures.
G-05 excludes positive differences with at most three rows.
The remaining partition has exactly four rows, as asserted.
\end{proof}
\paragraph{Effectivity and its limit.}
At the displayed degree, the finite coefficient and isotypic-rank
algorithms of R-05--R-06 can enumerate the partitions and find one
with a positive difference. The proof therefore bounds a finite
search, but its size is not polynomial. In the quartic case the
explicit family at degree $320112+s$ in R-12 is substantially
stronger than this uniform bound.


\begin{proposition}[G-11: Explicit rectangular multiplicity gaps]\label{G-11}
For $n\geq4$ define the following positive integers:
\[
\begin{gathered}
 b_n=\binom{n+3}{3},\qquad \beta_n=6b_n n^{30},\qquad
 T_n=\binom{b_n+\beta_n}{b_n}-1,\\
 E_n=4T_n(\beta_n!),\qquad u_n=2n^2-13,\\
 Q_n=(nE_n)^{u_n-1},\qquad K_n=u_n(Q_n-1)+1.
\end{gathered}
\]
For every integer $k\geq K_n$, put $d=E_n k$ and
$\lambda=(nd/4,nd/4,nd/4,nd/4)$. Then
\begin{equation}\label{G11:rectangle}
 m_\lambda(\operatorname{per}_n)>m_\lambda(\det_n).
\end{equation}
More generally, let $n>m\geq4$. The same conclusion, with
$\operatorname{per}_n$ replaced by
$\ell^{n-m}\operatorname{per}_m$, holds if
\begin{equation}\label{G11:padded-hypotheses}
 \dim\mathcal P_{m,4}+3>2n^2+2
 \quad\hbox{and}\quad n<4m/3.
\end{equation}
In particular it holds for $n=m+1$ and every $m\geq8$, and
for $n=\lfloor\rho m\rfloor$, all sufficiently large $m$,
and any fixed $1<\rho<\sqrt5/2$.
The degree, its progression, and the rectangular partition are
explicit. No useful computational size or polynomial algorithm
is asserted for these bounds.
\end{proposition}
\paragraph{Background used.}
We use the Kempf--Ness closed-orbit criterion in Losev (2006),
Theorem 4, p.~1; the stable-locus and quotient statements of
Mukai (2003), Definition 5.12, Proposition 5.15, Theorem 5.16,
and Corollary 5.17, pp.~165--167; and Derksen (2001),
Theorem 1.1, Proposition 1.2, and Example 2.1, printed p.~956.
The latter gives the displayed bound $\beta_n$ for generators
of the invariant ring of quaternary degree-$n$ forms.
The argument below proves the particular stability witnesses,
Veronese bound, and GCT deduction that those sources do not state.
\begin{proof}
Write $G=\mathrm{SL}_4(\mathbb C)$ and let $V_n$ be the space
of degree-$n$ quaternary forms. First we exhibit stable points
in the cones whose dimensions we compare. The form
\[
 f_n=x_0^n+x_1^n+x_2^n+x_3^n
\]
belongs to both $\mathcal P_{n,4}$ and $\mathcal D_{n,4}$.
Indeed factor $x_0^n+x_1^n$ and $x_2^n+x_3^n$ into binary
linear factors, place them on the diagonal and directed cycle
of an $n\times n$ pencil, and use its two perfect matchings.
For the determinant, multiply one cycle entry by the sign
$(-1)^{n-1}$.

Equip $V_n$ with an $\mathrm{SU}_4$-invariant Hermitian inner
product and use unit monomial vectors. Distinct monomials are
orthogonal. In the normalized sum of the four pure powers,
the mean exponent of each variable is $n/4$. Its moment map
therefore vanishes on the diagonal trace-zero Lie algebra.
It also vanishes on the off-diagonal directions: an operator
$x_j\partial/\partial x_i$, $i\ne j$, cannot carry one of
these pure powers to another when $n\geq3$. Thus the moment
map is zero. The Kempf--Ness criterion makes the $G$-orbit
closed. Its infinitesimal stabilizer is zero, since the
monomials $x_i^{n-1}x_j$ and $x_i^n$ in the derivative of
$f_n$ are distinct for $n\geq3$ and force all matrix entries
to vanish. Its stabilizer is consequently finite. This is a
stable point in each cone.

For the padded case put $a=n-m>0$ and suppose $3a<m$.
Let $e_0$ be the unit monomial $x_0^n$ and $e_i$ the unit
monomial proportional to $x_0^a x_i^m$, $1\leq i\leq3$.
Set
\[
 \theta_0=\frac{m-3a}{4m},\qquad
 \theta_1=\theta_2=\theta_3=\frac{n}{4m},\qquad
 F=\sum_{i=0}^3\sqrt{\theta_i}\,e_i.
\]
All four numbers are positive and sum to one. The mean
exponent of $x_0$ is $n\theta_0+3a\theta_1=n/4$, and
that of $x_i$ is $m\theta_i=n/4$. No two supported monomials
differ by a root $e_i-e_j$, because their differences are
multiples of $m\geq4$. The same moment-map calculation
therefore gives zero. We have $F=x_0^a q$, where
$q=\sum_{i=0}^3 c_i x_i^m$ and all $c_i\ne0$.
The preceding binary-factor construction puts $q$ in
$\mathcal P_{m,4}$, so $F\in\mathcal Z_{m,n}$.

Here is a direct finite-stabilizer check. For a linear vector
field $D_A=\sum_i L_i\partial/\partial x_i$, the equation
$D_A(F)=0$, after division by $x_0^{a-1}$, is
\[
 aL_0q+x_0D_Aq=0.
\]
Reduction modulo $x_0$ shows $L_0=d_0x_0$, because $a\ne0$
and $q(0,x_1,x_2,x_3)\ne0$. It follows that
$D_Aq=-ad_0q$. The distinct mixed monomials
$x_i^{m-1}x_j$ force every off-diagonal entry of $A$ to
vanish. The $x_0^m$ coefficient gives $(m+a)d_0=0$,
and the remaining pure powers give every other diagonal
entry zero. Thus the Lie stabilizer is zero and $F$ is
stable by Kempf--Ness.

For every reduced irreducible cone $X$ just considered,
the stable locus is a nonempty open subset and its quotient
fibers are the $15$-dimensional $G$-orbits (Mukai, loc. cit.).
Hence the dimension-of-fibers theorem yields
\begin{equation}\label{G11:quotient-dimension}
 \dim\mathbb C[X]^G=\dim X-15.
\end{equation}
In particular, writing $A_D=\mathbb C[\mathcal D_{n,4}]^G$
and $A_X=\mathbb C[X]^G$, G-06, R-15, and the first
hypothesis in \eqref{G11:padded-hypotheses} imply
\[
 \dim A_D\leq u_n,\qquad \dim A_X\geq u_n+1.
\]
The invariant functor is exact in characteristic zero.
Thus both invariant rings are graded quotients of
$\mathbb C[V_n]^G$, including when $X$ is padded; see also
Mukai, Proposition 5.10 and its proof, pp.~164--165.
Derksen's bound, specialized to $q=4$ and symmetric degree
$n$, says that this ambient invariant ring is generated in
degrees at most $\beta_n=6b_n n^{30}$. The dual convention
for forms makes no difference: the contragredient
transforms the action by an automorphism of $\mathrm{SL}_4$.
We may take at most $T_n$ homogeneous generators, because
the vector space of all coefficient polynomials in $b_n$
variables of positive degree at most $\beta_n$ has dimension
$T_n$.

We check explicitly that the $E_n$th Veronese of each ring
is generated in its first degree. More generally, suppose
an algebra has $r\leq T_n$ homogeneous generators with
positive degrees $d_i\leq\beta_n$, and put
$B_0=\beta_n!$. In a monomial in these generators of total
degree $kE_n$, remove as many pure blocks of
$B_0/d_i$ copies of generator $i$ as possible. Every removed
block has degree $B_0$. The remaining exponents are less
than $B_0/d_i$, so the residual degree is less than
$rB_0\leq T_n B_0$. It is an integer multiple $hB_0$,
because the original and removed degrees are such multiples;
thus $0\leq h<T_n$. There are $4kT_n-h$ pure blocks.
Combine the residual monomial with $4T_n-h$ of these blocks
to form a monomial of degree $E_n=4T_nB_0$, then divide
the other $4(k-1)T_n$ blocks into groups of $4T_n$.
This factors the original monomial into $k$ monomials of
degree $E_n$. Applying this to spanning monomials proves
the claim, also after passing to a quotient.

Consequently $A_D^{(E_n)}$ and $A_X^{(E_n)}$ are standard
graded domains, with dimensions unchanged. Set
$s=\dim\operatorname{Proj}A_D^{(E_n)}\leq u_n-1$.
This projective dimension is positive: block padding embeds
$\mathcal D_{4,4}$, whose dimension is $34$ by G-07,
into $\mathcal D_{n,4}$. By \eqref{G11:quotient-dimension}
the latter invariant ring has dimension at least $19$,
so $s\geq18$.
Choose a basis of its degree-one piece represented by
homogeneous invariant polynomials of degree $E_n$ on $V_n$.
Their composition with the determinant-pencil map is a
collection of homogeneous polynomials of degree $nE_n$.
They give a dominant rational map from
$\mathbb P^{4n^2-1}$ to $\operatorname{Proj}A_D^{(E_n)}$.
Dominance follows from dominance of the pencil map and
surjectivity of the affine quotient (Mukai, Theorem 5.9);
the Veronese morphism has this projective image. G-10's
base-locus-aware image-degree argument gives
\[
 \deg\operatorname{Proj}A_D^{(E_n)}
 \leq(nE_n)^s\leq Q_n.
\]
Chardin's bound in G-10 therefore gives, for every $k\geq1$,
\[
 \dim(A_D)_{E_nk}
 \leq Q_n\binom{k+u_n-1}{u_n-1}.
\]
Linear Noether normalization of the standard graded ring
$A_X^{(E_n)}$, whose dimension is at least $u_n+1$, gives
\[
 \dim(A_X)_{E_nk}\geq\binom{k+u_n}{u_n}.
\]
The ratio of the binomials is $(k+u_n)/u_n>Q_n$ when
$k\geq K_n$. Thus the first invariant dimension is strictly
smaller than the second.

Finally, in coefficient degree $d$, every polynomial
$\mathrm{GL}_4$ weight has size $nd$. The only irreducible
polynomial representation with trivial restriction to
$\mathrm{SL}_4$ is a determinant power; it has shape
$(nd/4)^4$, when $4\mid nd$, and is one-dimensional.
Here $E_n$ is divisible by four. The invariant dimension
difference is therefore exactly the difference of the two
multiplicities for the displayed rectangular partition.
The highest-weight restriction argument of G-05 and R-14
identifies these with the full orbit-closure multiplicities.
This proves \eqref{G11:rectangle} and its padded version.
R-14 and R-16 supply the stated particular ranges; all
satisfy $n<4m/3$.
\end{proof}
\paragraph{Relation to the rectangular bridge.}
The Haar projector onto $\mathrm{SU}_4$ invariants isolates
precisely these rectangular highest weights. Equivalently,
the corresponding character integral is the dimension of
this invariant graded piece. The proof makes that positive
rank exceed the determinant rank; it does not replace an
orbit-closure multiplicity by an ambient Kronecker
coefficient. The quartic formula G-07/R-12 is much sharper
where it applies. These deductions make no literature-priority
claim and are independent of RH and either proposed GRH input.

\begin{proposition}[G-12: A padding exclusion for the rectangular sector]\label{G-12}
Let $m\geq1$, $n>m$, and $n\geq4m/3$. For every degree-$m$
form $f$, every coordinate-ring degree $d\geq1$ with
$4\mid nd$, and $\lambda=(nd/4)^4$, one has
\begin{equation}\label{G12:exclusion}
 m_\lambda(\ell^{n-m}f)\leq m_\lambda(\det_n).
\end{equation}
If $n>4m/3$, the multiplicity on the left is zero.
The equality case $n=4m/3$ is included in the inequality.
Thus the rectangular gaps of G-11 cannot extend to or beyond
this padding boundary, even though nonrectangular four-row
questions remain.
\end{proposition}
\begin{proof}
As in G-05, restrict substituted forms to four variables.
They have shape $L^a q$, where $a=n-m$ and $q$ has degree
$m$. If $L=0$, the form is zero. Otherwise a change of
coordinates takes $L$ to $x_0$; it suffices to use a
$\mathrm{GL}_4$ change, since the determinant cone and each
homogeneous invariant ideal are stable under this group.
Act on the variables by the one-parameter subgroup
\[
 \eta(t)=\operatorname{diag}(t^3,t^{-1},t^{-1},t^{-1})
 \quad\hbox{of }\mathrm{SL}_4.
\]
A monomial in $x_0^a q$ with $x_0$ exponent $a+j$ has
weight $3a-m+4j$. If $3a>m$, all these weights are positive,
so $\eta(t)x_0^a q\longrightarrow0$ as $t\longrightarrow0$.
Every homogeneous positive-degree invariant therefore
vanishes on every such padded form and on their orbit
closure. Exactness of invariants identifies its invariant
coordinate ring with $\mathbb C$, whose positive-degree pieces vanish,
which proves the zero assertion.

If $3a=m$, the same limit instead equals
\[
 F_0=x_0^a q(0,x_1,x_2,x_3).
\]
Every ternary degree-$m$ form belongs to $\mathcal D_{m,3}$
by G-05 (including the elementary degrees one and two).
Embedding it and adding $a$ diagonal copies of $x_0$
shows $F_0\in\mathcal D_{n,4}$; the argument also holds for
limits because the block-padding map is polynomial.
Let $P$ be any homogeneous $\mathrm{SL}_4$ invariant
vanishing on $\mathcal D_{n,4}$. Since $P$ is constant
along $\eta$, continuity gives
$P(x_0^a q)=P(F_0)=0$. Thus the determinant ideal intersected
with the ambient invariant ring is contained in the ideal
of the restricted padded orbit closure intersected with
that ring. Both invariant coordinate rings are quotients
of the same ambient invariant ring, by exactness of
invariants. This containment gives a surjection from the
determinant invariant coordinate ring to the padded one,
and hence the degreewise dimension inequality.
Invariant dimensions are precisely the rectangular
multiplicities, as in G-11; highest-weight restriction
transfers the inequality to the full orbit closures.
\end{proof}
\paragraph{Source and scope.}
The strict inequality is also the rectangular specialization
of the usual padding weight constraint, but the proof above
gives its one-parameter-subgroup form and handles the
boundary by the ternary determinant theorem. This is an
exclusion for this specified rectangular sector. It does
not compare all four-row multiplicities or arbitrary
Kronecker coefficients.

\begin{proposition}[G-13: An explicit finite bound for the first rectangular gap]\label{G-13}
For $n\geq4$ put
\[
 \beta=6\binom{n+3}{3}n^{30},\qquad u=2n^2-13,\qquad
 C=(n\beta^2)^u,
\]
and define
\begin{equation}\label{G13:first-bound}
 s_*=(u+1)(C-1)+1,\qquad d_* = \beta s_*.
\end{equation}
There is an integer $d$ with $1\leq d\leq d_*$ and $4\mid nd$
such that, for $\lambda=(nd/4)^4$,
\[
 m_\lambda(\operatorname{per}_n)>m_\lambda(\det_n).
\]
The same conclusion holds with $\operatorname{per}_n$ replaced by
$\ell^{n-m}\operatorname{per}_m$ whenever $n>m\geq4$ and
\[
 \dim\mathcal P_{m,4}+3>2n^2+2,
 \qquad n<4m/3.
\]
In particular it applies to $n=m+1$, $m\geq8$, and, by R-18,
to $n=\lfloor\rho m\rfloor$ for all sufficiently large $m$
and each fixed $1<\rho<\sqrt{3/2}$.

The bound satisfies $\log d_*=O(n^2\log n)$. This is an
existence bound within a finite interval. It does not specify
the individual successful degree or replace G-11's assertion
at every degree of its displayed progression.
\end{proposition}

\paragraph{Background and scope.}
We use the stable points, exact invariant quotients, and quotient
dimensions established in G-11, together with Derksen (2001),
Theorem~1.1 and Example~2.1, printed p.~956. We also use the
projective image-degree argument of G-10 and Chardin's
unnumbered main theorem (1989), printed p.~306. The argument
below compares a weighted degree filtration with the ordinary
degree filtration in an affine generator embedding; it needs
no factorial Veronese.

\begin{proof}
Let $G=\mathrm{SL}_4(\mathbb C)$ and let $V_n$ be the space
of quaternary forms of degree $n$. Write
\[
 A_D=\mathbb C[\mathcal D_{n,4}]^G,
 \qquad A_X=\mathbb C[X]^G,
\]
where $X$ is $\mathcal P_{n,4}$ or the padded cone in the
statement. The stable-point and dimension arguments of G-11
give
\begin{equation}\label{G13:dims}
 t:=\dim A_D\leq u,\qquad \dim A_X\geq u+1.
\end{equation}
These are positively graded domains, with degree-zero part
$\mathbb C$, and both are graded quotients of
$\mathbb C[V_n]^G$. Choose homogeneous generators
$g_1,\ldots,g_r$ of the latter ring with positive degrees
$e_i\leq\beta$. Their images generate both quotient rings.

First bound the cumulative graded dimension of $A_D$. The
images of the generators give a closed affine embedding
\[
 Y_D=\operatorname{Spec} A_D\lhook\joinrel\longrightarrow
 \mathbb A^r.
\]
Compose this embedding with the determinant-pencil map and
the affine quotient. The resulting map from
$\mathbb A^{4n^2}$ has coordinates
$g_i(\det L)$, homogeneous of degrees $ne_i\leq n\beta$
in the pencil parameters. Its image is dense in $Y_D$.
To homogenize it, adjoin a source coordinate $z$ and use
the degree-$n\beta$ coordinate polynomials
\[
 z^{n\beta},\qquad
 z^{n\beta-ne_i}g_i(\det L),\quad 1\leq i\leq r.
\]
They define a rational map from $\mathbb P^{4n^2}$ whose
image closure is the ordinary projective closure
$\overline{Y_D}\subseteq\mathbb P^r$. It has dimension
$t$. G-10's image-degree bound, allowing a base locus, gives
\begin{equation}\label{G13:affine-degree}
 \deg\overline{Y_D}\leq(n\beta)^t\leq(n\beta)^u.
\end{equation}

For $d\geq0$, let
\[
 H_D^{\leq}(d)=\sum_{j=0}^d\dim(A_D)_j.
\]
Every element of $(A_D)_j$ has a preimage homogeneous of
weighted degree $j$ in $\mathbb C[T_1,\ldots,T_r]$ with
$\deg T_i=e_i$. Since every $e_i\geq1$, each monomial in
that preimage has ordinary degree at most $j$. It follows
that the direct sum of the graded pieces through $d$ is
contained in the space of functions on $Y_D$ represented
by ordinary polynomials of degree at most $d$.
This is an inclusion of filtered vector spaces, not an
identification of the two gradings.

The latter ordinary filtered dimension equals the
degree-$d$ Hilbert function of $\overline{Y_D}$:
homogenization with its first coordinate, followed by
dehomogenization on the dense affine chart, identifies the
two quotients. The projective closure is integral over
$\mathbb C$. Also $t\geq19$ by the fixed block-padding and
stable-quotient observation in G-11, so the positive-dimension
hypothesis in Chardin's stated theorem holds. That theorem
and \eqref{G13:affine-degree} yield
\begin{equation}\label{G13:upper}
 H_D^{\leq}(d)
 \leq(n\beta)^t\binom{d+t}{t}
 \leq(n\beta)^u\binom{d+u}{u}.
\end{equation}

For the other ring, a maximal algebraically independent
subset of the images of $g_1,\ldots,g_r$ has size
$\dim A_X$. Indeed those images generate its fraction
field, and maximality makes every remaining generator
algebraic over the chosen subset. By \eqref{G13:dims},
choose $u+1$ algebraically independent images. Each has
degree at most $\beta$. Their monomials of ordinary total
degree at most $s$ are linearly independent and have
graded degree at most $\beta s$. Consequently
\begin{equation}\label{G13:lower}
 H_X^{\leq}(\beta s)
 :=\sum_{j=0}^{\beta s}\dim(A_X)_j
 \geq\binom{s+u+1}{u+1}.
\end{equation}

For every nonnegative integer $s$, the exact binomial ratio
satisfies
\begin{align*}
 \frac{\binom{s+u+1}{u+1}}{\binom{\beta s+u}{u}}
 &=\frac{s+u+1}{u+1}
     \prod_{i=1}^u\frac{s+i}{\beta s+i}\\
 &\geq\frac{s+u+1}{(u+1)\beta^u}.
\end{align*}
The last inequality follows term by term from
$\beta(s+i)\geq\beta s+i$. At $s=s_*$ one has
\[
 s_*+u+1=(u+1)C+1,
\]
so the final lower bound is strictly greater than
$C/\beta^u=(n\beta)^u$. Equations
\eqref{G13:upper}--\eqref{G13:lower} therefore give
\[
 H_X^{\leq}(d_*)>H_D^{\leq}(d_*).
\]
The degree-zero dimensions are both one. Hence for some
$1\leq d\leq d_*$ the individual graded pieces satisfy
$\dim(A_X)_d>\dim(A_D)_d$.

As in G-11, the only polynomial $\mathrm{GL}_4$
representation of weight size $nd$ that is trivial on
$\mathrm{SL}_4$ is the one-dimensional determinant power
of shape $(nd/4)^4$. If $4\nmid nd$, both invariant
pieces are zero. Thus the successful degree necessarily
satisfies $4\mid nd$, and the invariant dimension gap is
the desired rectangular multiplicity gap. The first-four-
variable highest-weight argument in G-05/R-14 transfers
it to the full orbit closures.

Finally $\beta=O(n^{33})$ and $u=O(n^2)$, so the explicit
formula \eqref{G13:first-bound} gives
$\log d_*=O(n^2\log n)$. The particular padded ranges
follow from R-14/R-16 and R-18, respectively; their
dimension gaps and $n<4m/3$ are the same hypotheses already
used to construct the stable padded point in G-11.
\end{proof}

\paragraph{Complexity lower-bound context.}
The multiplicity statements G-10--G-13 locate particular representation
sectors and degrees. They do not improve the established quadratic
lower bound $\overline{\operatorname{dc}}(\operatorname{per}_m)
\geq m^2/2$ of Landsberg, Manivel, and Ressayre (2013);
the inspected April 2010 author version states it as Theorem~1.0.1,
p.~1 (G-B16). No claim about the strongest currently proposed lower
bound is needed here. A noncontainment result by itself does not
specify the four-row or rectangular multiplicity differences proved
in the present restricted parameter ranges.
